\shorttitle{Двомасштабне стохастичне моделювання з агрегованим міжспільнотним зв'язком}
\shortauthors{Y. Kuryliak and M.T.M. Emmerich}

\title[mode=title]{Двомасштабне стохастичне моделювання епідемій у мережах спільнот з агрегованим міжспільнотним зв'язком}

\author[1]{Yulian Kuryliak}[orcid=0000-0002-6600-2637]
\cormark[1]
\ead{yulian.a.kuryliak@lpnu.ua}
\affiliation[1]{organization={Інститут комп'ютерних наук та інформаційних технологій, Національний університет ``Львівська політехніка''},
            addressline={вул. Степана Бандери, 12},
            city={Львів},
            postcode={79000},
            country={Україна}}

\author[2]{Michael T. M. Emmerich}[orcid=0000-0002-7342-2090]
\ead{michael.t.m.emmerich@jyu.fi}
\affiliation[2]{organization={Факультет інформаційних технологій, Університет Ювяскюля},
            addressline={P.O. Box 35 (Agora)},
            city={Ювяскюля},
            postcode={FI-40014},
            country={Фінляндія}}

\cortext[1]{Автор для листування}

\begin{abstract}
У багатомасштабному стохастичному моделюванні детальна динаміка підсистем часто поєднується через агреговані інтерфейси. Ми досліджуємо цю задачу для поширення епідемій у мережах спільнот, де контакти всередині спільнот подано явно, а зв'язки між ними доступні лише в агрегованій формі. Кожна спільнота моделюється рушієм Gillespie, точним умовно на отриманих імпортованих подіях; макрорівень генерує імпорти за агрегованими залежними від стану інтенсивностями, а в базовій реалізації розміщує їх рівномірно серед сприйнятливих вузлів спільноти-одержувача. Ми повністю специфікуємо метод, зокрема його сіткову синхронізацію, і показуємо, що ідеальна версія з неперервним часом точно моделює сурогатну мережу, у якій кожен пучок міжспільнотних ребер замінено рівномірним зв'язком «усі з усіма» з такою самою сумарною вагою. Це подання виявляє чотири джерела розбіжностей: втрату інформації про стан граничних вузлів, рівномірне розміщення імпортів, відсутність вичерпання ребер інтерфейсу та похибку синхронізації. Для випадково розміщених інтерфейсів у задачі з одиничним експортом ми також доводимо, що ідеальний агрегований годинник експорту стохастично випереджає повномережевий. Контрольовані експерименти на синтетичних мережах SI порівнюють метод із повномережевим моделюванням Gillespie у тестах з одиничним експортом і в 48 конфігураціях ланцюгів із десяти спільнот (50--1000 вузлів у спільноті, внутрішні довжини шляхів 1--6, одне або п'ять ребер інтерфейсу, по 400 зіставлених запусків). Коли розкид часів зараження всередині спільноти був меншим за чверть середнього часу переходу через одне ребро, розподіли часу прибуття до кінцевої спільноти були статистично нерозрізненними у 15 із 18 конфігурацій; за більших значень двомасштабний симулятор прибував раніше на величину, приблизно пропорційну цьому відношенню, яка сягала 40\%. Близькі часи не гарантували допустимих маршрутів: до 52\% агрегованих переходів відбувалися до зараження будь-якого справжнього вузла-моста.
\end{abstract}

\begin{highlights}
\item Точні локальні рушії Gillespie зв'язані через агреговані міжспільнотні інтенсивності
\item Ідеальна схема точно моделює сурогатну мережу з інтерфейсом «усі з усіма»
\item Для випадкових інтерфейсів агрегований годинник експорту стохастично випереджає повномережевий
\item Розбіжність упорядковується співвідношенням внутрішнього розкиду часу та часу переходу через ребро
\item Близькі часи прибуття можуть супроводжуватися фізично неможливими маршрутами передавання
\end{highlights}

\begin{keywords}
Стохастичне моделювання \sep Багатомасштабне моделювання \sep Алгоритм Gillespie \sep Агреговані зв'язки \sep Точність моделі \sep Мережі спільнот
\end{keywords}

\maketitle

%======================================================================
\clearpage
\input{readability_revision_parallel_black_20261008/sections_1_3_UK.tex}

\clearpage
